%% AUTO-GENERATED by docs/gen-tactics-ref.py from reference/TACTICS.md -- do not edit.
%% Regenerate:  python3 docs/gen-tactics-ref.py   (after a prover load refreshes TACTICS.md)

\DeclareRobustCommand{\tactkind}[1]{{\footnotesize\textsc{#1}}}

Every tactic is tagged with a \emph{kind} grounded in the \texttt{dg-apply-rule!} tag it emits, so the trusted base is legible: \texttt{rule} is a single primitive kernel inference (the fixed trusted base); \texttt{oracle} is a trusted decision procedure run as a black box; \texttt{composite} only chains kernel rules, adding no new inference; \texttt{meta} performs no deduction (session, search, navigation).  A proof's trust surface is exactly its \texttt{rule} steps plus whichever \texttt{oracle}s and asserted premises it cites.

The same reference is available at the REPL with \texttt{(tactics 'name)} and in the browser file \texttt{reference/TACTICS.md}; all three are generated from one registry, so they never disagree.

An entry's \emph{Uses} line names every kernel operation the command's own code can cause to be recorded, those recorded by the procedures it calls included. It is the bound on what the command can do, whether or not any proof in the library exercises it, and it is measured by reading the source rather than declared alongside the command. One further set is common to every command and is stated once instead of in every entry: the hook that runs after any command which changed the proof, to close the definedness sequents an instantiation left owed. Section~\ref{sec:command-uses} tabulates the same figures for every command at once.

\subsection*{Starting \& finishing a proof}

\par\smallskip{\raggedright\noindent\texttt{(sp goal)}\quad\tactkind{meta}\index{sp@\texttt{sp} (tactic)}\par}\nobreak
\noindent Start a proof of `goal' (a ``string'', a raw S-expr, or a wff); clears the script.\par
\noindent\textit{Uses:} no kernel operation: it records no inference.\par
\noindent\textit{When useful:} starting a new proof\par

Begin a new proof: state what you want to prove.  The goal may be a ``string'' in surface syntax, a raw S-expr, or an already-built wff (e.g. (wff ``...'') or (make-wff '(...)));  sp coerces a string/S-expr for you.  This becomes your one open goal and clears any previous script.  (Technically: set the proof goal.)

\par\smallskip{\raggedright\noindent\texttt{(qed 'name)}\quad\tactkind{meta}\index{qed@\texttt{qed} (tactic)}\par}\nobreak
\noindent Install the finished proof as theorem `name' and save its replayable script.\par
\noindent\textit{Uses:} no kernel operation: it records no inference.\par
\noindent\textit{When useful:} no open goals remain -- record the result\par

Finish: once there are no open goals, this records the result as a named theorem you can cite later, and saves the replayable script.  (Technically: install-theorem! plus script capture; reports the asserted facts the proof still rests on, `proven modulo \{...\}'.)

\par\smallskip{\raggedright\noindent\texttt{(save-proof 'name)}\quad\tactkind{meta}\index{save-proof@\texttt{save-proof} (tactic)}\par}\nobreak
\noindent Save the current script under a name without finishing.\par
\noindent\textit{Uses:} no kernel operation: it records no inference.\par

Save the current (possibly unfinished) proof script under a name, to resume or replay later.  (Technically: snapshot the script without installing a theorem.)

\par\smallskip{\raggedright\noindent\texttt{(replay-proof 'name [subst])}\quad\tactkind{meta}\index{replay-proof@\texttt{replay-proof} (tactic)}\par}\nobreak
\noindent Re-run a saved script on the current goal, optionally renaming free vars.\par
\noindent\textit{Uses:} every one of the 64 kernel operations.\par

The optional subst is an alist ((old . new) ...) applied to every command argument before replay -- this is how one proof is reused at fresh eigenvariables.


\subsection*{Decomposing the goal}

\par\smallskip{\raggedright\noindent\texttt{(di)}\quad\tactkind{rule}\index{di@\texttt{di} (tactic)}\par}\nobreak
\noindent Direct inference: split an AND goal, move an IMPLIES antecedent into the assumptions, or introduce a fresh eigenvariable for a leading FORALL.\par
\noindent\textit{Uses:} \texttt{and-intro} \texttt{forall-intro} \texttt{iff-intro} \texttt{implies-intro} \texttt{not-intro}.\par
\noindent\textit{When useful:} the goal head is AND / IMPLIES / FORALL -- decompose before deciding\par

Break the goal into its pieces.  A conjunction `A and B' splits into two goals, A and B.  An implication `P implies Q' assumes P (P becomes a hypothesis you may use) and leaves you to prove Q.  A universal `for all x, ...' fixes an arbitrary x (a new constant standing for `any x') and asks for the statement at that x.  Apply it repeatedly until the goal is a single atomic statement.  (Technically: the goal-side introduction rules for AND / IMPLIES / FORALL; the constant introduced for a FORALL is an eigenvariable.  Pairs with mac, which unfolds a definition in the goal.)

\par\smallskip{\raggedright\noindent\texttt{(pbc)}\quad\tactkind{rule}\index{pbc@\texttt{pbc} (tactic)}\par}\nobreak
\noindent Proof by contradiction: assume the goal's negation, prove FALSITY.\par
\noindent\textit{Uses:} \texttt{proof-by-contradiction}.\par
\noindent\textit{When useful:} a direct argument stalls and assuming the negation gives something concrete to work with\par

Argue by contradiction: assume the goal is false and derive an absurdity.  (Technically: classical reductio -- adds the negation of the goal as a hypothesis and changes the goal to FALSITY.)

\par\smallskip{\raggedright\noindent\texttt{(oi-l)}\quad\tactkind{rule}\index{oi-l@\texttt{oi-l} (tactic)}\par}\nobreak
\noindent OR-intro left: reduce an (OR a b) goal to a.\par
\noindent\textit{Uses:} \texttt{or-intro-left}.\par
\noindent\textit{When useful:} the goal is (OR a b) and the LEFT alternative is provable\par

Prove a disjunction `A or B' by proving the LEFT alternative, A.  (Technically: or-introduction, left.)

\par\smallskip{\raggedright\noindent\texttt{(oi-r)}\quad\tactkind{rule}\index{oi-r@\texttt{oi-r} (tactic)}\par}\nobreak
\noindent OR-intro right: reduce an (OR a b) goal to b.\par
\noindent\textit{Uses:} \texttt{or-intro-right}.\par
\noindent\textit{When useful:} the goal is (OR a b) and the RIGHT alternative is provable\par

Prove a disjunction `A or B' by proving the RIGHT alternative, B.  (Technically: or-introduction, right.)

\par\smallskip{\raggedright\noindent\texttt{(ew term)}\quad\tactkind{rule}\index{ew@\texttt{ew} (tactic)}\par}\nobreak
\noindent Existential witness: discharge a FORSOME goal by supplying the witness term.\par
\noindent\textit{Uses:} \texttt{forsome-intro}.\par
\noindent\textit{When useful:} the goal is FORSOME and you have a witness term in mind\par

Prove `there exists an x with property P' by exhibiting a specific witness: you supply the term, and the goal becomes `P holds of that term'.  The everyday `take x = ...' step.  (Technically: existential introduction for a FORSOME goal.)

\par\smallskip{\raggedright\noindent\texttt{(ci)}\quad\tactkind{rule}\index{ci@\texttt{ci} (tactic)}\par}\nobreak
\noindent Cartesian intro: prove a CARTESIAN-product membership component-wise.\par
\noindent\textit{Uses:} \texttt{cartesian-intro}.\par
\noindent\textit{When useful:} the goal is membership in a CARTESIAN product\par

Prove that a pair (or tuple) belongs to a Cartesian product A x B by proving each coordinate lies in its factor.  (Technically: Cartesian-membership introduction, component-wise.)

\par\smallskip{\raggedright\noindent\texttt{(ti)}\quad\tactkind{rule}\index{ti@\texttt{ti} (tactic)}\par}\nobreak
\noindent Tuple intro: prove a tuple/LIST membership component-wise.\par
\noindent\textit{Uses:} \texttt{tuples-intro}.\par
\noindent\textit{When useful:} the goal is membership in a set of tuples / LIST\par

Prove that a list belongs to the set of tuples over A by proving each entry lies in A.  (Technically: tuple-membership introduction, component-wise.)

\par\smallskip{\raggedright\noindent\texttt{(ii)}\quad\tactkind{rule}\index{ii@\texttt{ii} (tactic)}\par}\nobreak
\noindent Intersection intro: prove (IN x (INTERSECTION ...)) for each branch.\par
\noindent\textit{Uses:} \texttt{intersection-intro}.\par
\noindent\textit{When useful:} the goal is membership in an INTERSECTION\par

Prove that x lies in an intersection of sets by proving it lies in each of them.  (Technically: intersection-membership introduction.)

\par\smallskip{\raggedright\noindent\texttt{(ui k)}\quad\tactkind{rule}\index{ui@\texttt{ui} (tactic)}\par}\nobreak
\noindent Union intro: reduce a goal (IN x (UNION ...)) to membership of x in the k-th set (1-based).\par
\noindent\textit{Uses:} \texttt{union-intro}.\par
\noindent\textit{When useful:} the goal is membership in a UNION and you know which branch\par

Prove that x lies in a union of sets by proving it lies in the k-th one (you choose which).  (Technically: union-membership introduction at branch k.)

\par\smallskip{\raggedright\noindent\texttt{(ni)}\quad\tactkind{rule}\index{ni@\texttt{ni} (tactic)}\par}\nobreak
\noindent Natural-number induction on the goal's leading FORALL over NN.\par
\noindent\textit{Uses:} \texttt{nn-induction}.\par
\noindent\textit{When useful:} the goal is `forall n in NN, P(n)'\par

Prove `for all natural numbers n, P(n)' by induction: show P(0), and that P(n) implies P(n+1).  (Technically: natural-number induction on the goal's leading FORALL over NN.)


\subsection*{Using a hypothesis}

\par\smallskip{\raggedright\noindent\texttt{(ai hyp)}\quad\tactkind{rule}\index{ai@\texttt{ai} (tactic)}\par}\nobreak
\noindent Antecedent inference: decompose a cited assumption -- AND-split, OR-into-cases, or FORSOME-elimination to a fresh eigenvariable.\par
\noindent\textit{Uses:} \texttt{and-elim} \texttt{forsome-elim} \texttt{iff-elim} \texttt{not-elim} \texttt{or-elim}.\par
\noindent\textit{When useful:} a hypothesis is an AND / OR / FORSOME to break apart\par

Break a hypothesis into its pieces -- the mirror image of di, but on the assumptions instead of the goal.  From a hypothesis `A and B' you get both A and B; from `A or B' you split into two cases (prove the goal in each); from `there exists x, P(x)' you get a fresh name for such an x together with P(x).  Cite the hypothesis by its number in the display, its formula, or a ``string''.  (Technically: the hypothesis-side elimination rules for AND / OR / FORSOME; the fresh name is an eigenvariable.)

\par\smallskip{\raggedright\noindent\texttt{(ass)}\quad\tactkind{rule}\index{ass@\texttt{ass} (tactic)}\par}\nobreak
\noindent Close the goal by an assumption alpha-equivalent to it.\par
\noindent\textit{Uses:} \texttt{assumption}.\par
\noindent\textit{When useful:} the goal already appears (up to bound-var renaming) among the hypotheses\par

Finish the goal because it is already one of your hypotheses -- the `that is exactly what we assumed' step.  Renaming of dummy/bound variables doesn't matter (`there exists x, P(x)' closes `there exists y, P(y)').  (Technically: the assumption rule, matching up to alpha-equivalence.)

\par\smallskip{\raggedright\noindent\texttt{(inst forall-hyp term)}\quad\tactkind{rule}\index{inst@\texttt{inst} (tactic)}\par}\nobreak
\noindent Instantiate a universally-quantified assumption at `term', adding the instance to context.\par
\noindent\textit{Uses:} \texttt{forall-elim}.\par
\noindent\textit{When useful:} you have a forall-hypothesis and a specific value to use it at\par

Use a `for all x, ...' hypothesis at a particular value: you supply the term, and the statement with x replaced by that term is added to your hypotheses.  (Technically: universal instantiation of an assumption.)

\par\smallskip{\raggedright\noindent\texttt{(inst+ forall-hyp term)}\quad\tactkind{composite}\index{inst+@\texttt{inst+} (tactic)}\par}\nobreak
\noindent Instantiate an in-context universal at `term', then forward-detach any guards whose antecedents are in context.\par
\noindent\textit{Uses:} \texttt{arith-forsome} \texttt{arith-ground} \texttt{arith-simplify} \texttt{cut} \texttt{detach} \texttt{forall-elim}.\par
\noindent\textit{When useful:} a GUARDED forall-hyp whose guards are dischargeable from context (e.g. the metric laws post-grind)\par

inst followed by detach: from `forall x. (x in S) => P(x)' and `t in S' already known, this lands `P(t)' directly (peeling nested guards level by level), instead of leaving the guarded implication for you to detach by hand.  The hypothesis-side analogue of `fact' (which assembles a THEOREM); this assembles an in-context UNIVERSAL.  scout's inst lane emits these to close witness-needing goals like the metric laws.  (Technically: pi-instantiate! then pi-detach! while the consequent stays a guard with an in-context antecedent.)

\par\smallskip{\raggedright\noindent\texttt{(mp)}\quad\tactkind{composite}\index{mp@\texttt{mp} (tactic)}\par}\nobreak
\noindent Close a goal that is an INSTANCE of a universal already in your context.  No arguments: the term is DERIVED by matching, not supplied.\par
\noindent\textit{Uses:} \texttt{arith-forsome} \texttt{arith-ground} \texttt{arith-simplify} \texttt{assumption} \texttt{cut} \texttt{detach} \texttt{forall-elim}.\par
\noindent\textit{When useful:} the goal IS an instance of a universal you already have -- no term to supply, it is derived by matching\par

The syllogism, with nothing to fill in.  From `forall([thing in human], thing in mortal)' together with `socrates in human', it closes `socrates in mortal'.  Where inst+ makes you name the term, mp works it out: it matches the universal's conclusion against the goal, which either determines the bound variable or fails outright, and then checks that the universal's guard at that term is something you already have.  Nothing is searched for and nothing is ranked -- if the goal is an instance there is exactly one term it can be.  It DECLINES, naming what it found, when no universal in the context has the goal as an instance, and when more than one does: a closer that silently picks between two universals is one you cannot predict.  One bound variable in this version; for `forall([a in A, b in B], ...)' use inst+.  (Technically: composite -- inst+ at the derived term, then ass.  It adds no kernel rule and no debt; the bill is what typing the two steps by hand would bill, which is nothing.)

\par\smallskip{\raggedright\noindent\texttt{(grind-and-mp)}\quad\tactkind{composite}\index{grind-and-mp@\texttt{grind-and-mp} (tactic)}\par}\nobreak
\noindent Tidy the goal with grind, then close it by modus ponens -- the one-button form of the syllogism, for a sentence typed exactly as it reads.\par
\noindent\textit{Uses:} \texttt{and-elim} \texttt{and-intro} \texttt{arith-forsome} \texttt{arith-ground} \texttt{arith-simplify} \texttt{assumption} \texttt{cut} \texttt{detach} \texttt{forall-elim} \texttt{forall-intro} \texttt{forsome-elim} \texttt{iff-elim} \texttt{iff-intro} \texttt{implies-intro} \texttt{macete-hyp} \texttt{not-elim} \texttt{not-intro} \texttt{or-elim}.\par
\noindent\textit{When useful:} a sentence typed as it reads -- `all men are mortal and socrates is a man implies socrates is mortal' -- closing in one move\par

Runs the two steps a syllogism actually takes on a sentence you have just typed.  `forall([thing in human], thing in mortal) and socrates in human implies socrates in mortal' is not yet a syllogism: it is an implication whose antecedent is a conjunction.  grind does the BOOKKEEPING -- split the AND, move the antecedents into the assumptions -- which leaves the goal `socrates in mortal' with the universal and `socrates in human' as hypotheses; mp then does the LOGIC (universal instantiation, then modus ponens).  They stay separate tactics because they are different kinds of step, and mp on its own is what to reach for once the goal is already tidy.  The recorded script keeps both steps, so the page reads as the two moves rather than as this wrapper.  (Technically: composite -- grind, quietly since on a tidy goal it legitimately declines, then mp.  Adds no kernel rule and no debt.)

\par\smallskip{\raggedright\noindent\texttt{(detach! impl)}\quad\tactkind{rule}\index{detach"!@\texttt{detach"!} (tactic)}\par}\nobreak
\noindent Forward modus ponens: from an in-context (IMPLIES A B) whose A is also in context, leave B in context.\par
\noindent\textit{Uses:} \texttt{detach}.\par
\noindent\textit{When useful:} you have both P and (P => Q) in context and want Q\par

If you have both `P' and `P implies Q' among your hypotheses, this adds `Q' to them.  It grows what you KNOW (the hypotheses) rather than changing the goal -- the forward counterpart of bc.  (Technically: forward modus ponens, the kernel rule pi-detach!; sound since P and P=>Q give Q.)

\par\smallskip{\raggedright\noindent\texttt{(fact 'thm term ...)}\quad\tactkind{composite}\index{fact@\texttt{fact} (tactic)}\par}\nobreak
\noindent Forward APPLICATION of a theorem: bring it in, instantiate its leading universals with the terms, and auto-detach every antecedent already in context, landing the consequent as a hypothesis.\par
\noindent\textit{Uses:} \texttt{arith-forsome} \texttt{arith-ground} \texttt{arith-simplify} \texttt{cut} \texttt{detach} \texttt{forall-elim} \texttt{theorem-assumption}.\par
\noindent\textit{When useful:} a library law `forall x. H(x) => P(x)' whose P you want landed as a hypothesis\par

Handles interleaved forall/implies (e.g. forall s. IS-X(s) => forall a. a in CARR(s) => P): consumes one term per FORALL, detaches each IMPLIES whose antecedent is in context.  The forward-assembly workhorse -- a law `forall x. H(x) => P(x)' becomes the usable fact P in one call, instead of ta + inst* + cut/backchain.  See theorem-library/module-zero-act.scm.

\par\smallskip{\raggedright\noindent\texttt{(ce hyp k)}\quad\tactkind{rule}\index{ce@\texttt{ce} (tactic)}\par}\nobreak
\noindent Cartesian elim: project the k-th component out of a CARTESIAN-membership assumption.\par
\noindent\textit{Uses:} \texttt{cartesian-elim}.\par
\noindent\textit{When useful:} a hypothesis is CARTESIAN-product membership\par

From a hypothesis that a tuple lies in a Cartesian product, extract that its k-th coordinate lies in the k-th factor.  (Technically: Cartesian-membership elimination, k-th projection.)

\par\smallskip{\raggedright\noindent\texttt{(te hyp k)}\quad\tactkind{rule}\index{te@\texttt{te} (tactic)}\par}\nobreak
\noindent Tuple elim: project the k-th component out of a tuple-membership assumption.\par
\noindent\textit{Uses:} \texttt{tuples-elim}.\par
\noindent\textit{When useful:} a hypothesis is tuple membership\par

From a hypothesis that something is a tuple over A, extract that its k-th entry lies in A.  (Technically: tuple-membership elimination, k-th projection.)

\par\smallskip{\raggedright\noindent\texttt{(ie hyp k)}\quad\tactkind{rule}\index{ie@\texttt{ie} (tactic)}\par}\nobreak
\noindent Intersection elim: extract the k-th branch of an INTERSECTION-membership assumption.\par
\noindent\textit{Uses:} \texttt{intersection-elim}.\par
\noindent\textit{When useful:} a hypothesis is INTERSECTION membership\par

From a hypothesis that x lies in an intersection, extract that x lies in the k-th set.  (Technically: intersection-membership elimination.)

\par\smallskip{\raggedright\noindent\texttt{(ue hyp)}\quad\tactkind{rule}\index{ue@\texttt{ue} (tactic)}\par}\nobreak
\noindent Union elim: split a cited (IN x (UNION ...)) membership assumption into one subgoal per set -- union-side case analysis, the dual of `ui'.\par
\noindent\textit{Uses:} \texttt{union-elim}.\par
\noindent\textit{When useful:} a hypothesis is UNION membership -- case-split on it\par

From a hypothesis that x lies in a union, split into cases -- one for each set x might belong to -- and prove the goal in each.  (Technically: union-membership elimination, case analysis; the dual of ui.)

\par\smallskip{\raggedright\noindent\texttt{(cut formula)}\quad\tactkind{rule}\index{cut@\texttt{cut} (tactic)}\par}\nobreak
\noindent Cut: prove `formula' as a side subgoal, then continue the main goal with `formula' added as an assumption (Gentzen cut).\par
\noindent\textit{Uses:} \texttt{cut}.\par
\noindent\textit{When useful:} you want to prove a lemma on the spot, then use it\par

Introduce a lemma you prove on the spot.  You state a formula; it becomes a side goal to prove, and on the main line you may then use it as a hypothesis.  The standard `we first show that ..., and now using it ...' move.  (Technically: Gentzen's cut -- nothing is left assumed-but-unproved, the side goal discharges it.)

\par\smallskip{\raggedright\noindent\texttt{(wk hyp)}\quad\tactkind{rule}\index{wk@\texttt{wk} (tactic)}\par}\nobreak
\noindent Weaken: drop a cited assumption from the context to tidy the hypothesis list.  `hyp' may be a formula, a ``string'', or a 1-based assumption index.\par
\noindent\textit{Uses:} \texttt{weakening}.\par
\noindent\textit{When useful:} the hypothesis list is cluttered with something no longer needed\par

Discard a hypothesis you no longer need, to keep the assumption list readable.  (Technically: weakening -- removing a hypothesis is always sound.)

\par\smallskip{\raggedright\noindent\texttt{(keep hyp ...)}\index{keep@\texttt{keep} (tactic)}\par}\nobreak
\noindent Keep only the cited assumptions and drop every other one, as ONE recorded step.  Each `hyp' may be a formula, a ``string'', or a 1-based assumption index; they are recorded as formulas.\par
\noindent\textit{When useful:} the context is deep and only a few hypotheses matter for the next step (prop, ineq)\par

Prune the context down to what the next step needs -- what a driver does before `prop' or `ineq' in a deep context.  (Technically: one weakening per dropped hypothesis, on the same leaf, recorded once; the driver kit's dk-only! is this command.)


\subsection*{Rewriting}

\par\smallskip{\raggedright\noindent\texttt{(mac 'name)}\quad\tactkind{rule}\index{mac@\texttt{mac} (tactic)}\par}\nobreak
\noindent Rewrite the GOAL with an equivalence macete (unfold a definition, apply an iff/=/== law).  Fires only where the macete's side-conditions already hold in context.\par
\noindent\textit{Uses:} \texttt{cartesian-decompose} \texttt{macete} \texttt{tuple-equality-decompose}.\par
\noindent\textit{When useful:} the GOAL has a defined predicate to unfold / an identity to apply, side-conditions already met\par

Rewrite the goal using a definition or a known equivalence/equality (a `macete').  For instance, replace a defined predicate by what it stands for, or apply an identity.  It fires only where the law's side-conditions already hold in your hypotheses; where they don't, it leaves that spot untouched and looks deeper inside.  (Technically: goal-side rewriting by an equivalence/equality macete; all-or-nothing -- it does NOT spawn unmet side-conditions as goals.  Its hypothesis-side cousin is mac-h; the minor-premise-spawning variant is macm.)

\par\smallskip{\raggedright\noindent\texttt{(slot 'acc)}\index{slot@\texttt{slot} (tactic)}\par}\nobreak
\noindent Reduce a structure ACCESSOR to its projection: (CARR s) becomes (NTH 1 s).  The one door for accessor reductions -- use it, never mac, on an accessor name.\par
\noindent\textit{Uses:} \texttt{cartesian-decompose} \texttt{macete} \texttt{tuple-equality-decompose}.\par

Replace an accessor by the tuple position it stands for: (CARR s) is slot 1, so it becomes (NTH 1 s).  Pair it with nth-r to compute on a concrete structure -- (MUL ZZ-RING) reduces to bintimes.  It refuses anything that is not an accessor.  (Technically: fires the accessor's macete, which def-structure installs at declaration.  It exists so that ALL accessor reductions go through ONE procedure: the reduction is currently global and unconditional -- index k for every argument -- and making it structure-relative later, guarded by IS-X(s), should change this door and not its callers.  accessor-callsite-audit fails the suite if any file fires an accessor macete by name.)

\par\smallskip{\raggedright\noindent\texttt{(slot-h 'acc hyp)}\index{slot-h@\texttt{slot-h} (tactic)}\par}\nobreak
\noindent Reduce a structure ACCESSOR to its projection inside a cited HYPOTHESIS -- `slot', hypothesis-side.\par
\noindent\textit{Uses:} \texttt{macete-hyp}.\par

The hypothesis-side door of slot: replace (PTS RR-MS) by the tuple position it stands for inside an assumption you cite, instead of in the goal.  Before it existed, a proof needing that reduction in a hypothesis reached for mac-h with the accessor macete's own name (rr-ms@pts) -- firing an accessor macete BY NAME, which is what the accessor pin (accessor-callsite-audit) exists to prevent, and which kept the suite red.  It refuses anything that is not an accessor, and it ERRORS rather than guess when the hypothesis mentions two different instances of the same accessor: after the first rewrite the assumption is a different formula, so a silent half-rewrite would be the failure mode.  (Technically: slot's projection-macete lookup run over the cited assumption, then cmd-apply-macete-to-assumption.)

\par\smallskip{\raggedright\noindent\texttt{(macm 'name)}\quad\tactkind{rule}\index{macm@\texttt{macm} (tactic)}\par}\nobreak
\noindent Like mac, but a CONDITIONAL macete fires even when its side-conditions are not yet in context: each unmet condition is left as a new subgoal.\par
\noindent\textit{Uses:} \texttt{cartesian-decompose} \texttt{macete} \texttt{tuple-equality-decompose}.\par
\noindent\textit{When useful:} the GOAL has a CONDITIONAL identity/definition to apply whose side-conditions you are willing to leave as subgoals\par

The minor-premises variant of mac.  mac only rewrites where the law's side-conditions ALREADY hold; macm rewrites regardless, spawning each unmet condition as a fresh goal to discharge -- the goal-side mirror of what mac-h already does on hypotheses.  This is what makes conditional finite-sum lemmas (finsum-*, sum-*, ...) usable as goal rewrites instead of only forward via fact.  (Technically: goal-side rewriting by a conditional macete; it emits the SAME single `macete' kernel rule as mac -- the unmet conditions become that one step's minor-premise subgoals.  It introduces NO new inference rule.)

\par\smallskip{\raggedright\noindent\texttt{(mac-h 'name hyp)}\quad\tactkind{rule}\index{mac-h@\texttt{mac-h} (tactic)}\par}\nobreak
\noindent Rewrite a cited ASSUMPTION in place with an equivalence macete; any side-condition not already in context is spawned as a new subgoal.\par
\noindent\textit{Uses:} \texttt{macete-hyp}.\par
\noindent\textit{When useful:} a HYPOTHESIS has a definition to unfold / an equivalence to apply\par

Like mac, but rewrites inside one of your HYPOTHESES instead of the goal -- replacing it by an equivalent statement (e.g. unfolding a definition in a hypothesis).  If the rewrite law has a side-condition you haven't established, that side-condition becomes a new goal to prove.  (Technically: hypothesis-side rewriting by Leibniz substitution of equivalents; sound because the macete is a genuine equivalence under its side-conditions, which are spawned as goals.)

\par\smallskip{\raggedright\noindent\texttt{(mac-h*)}\quad\tactkind{composite}\index{mac-h*@\texttt{mac-h*} (tactic)}\par}\nobreak
\noindent Saturating mac-h: repeatedly unfold every defined predicate in the hypotheses and split the conjunctions they expose, until nothing is left folded.  One step in the trace instead of a dozen.\par
\noindent\textit{Uses:} \texttt{and-elim} \texttt{forsome-elim} \texttt{iff-elim} \texttt{macete-hyp} \texttt{not-elim} \texttt{or-elim}.\par
\noindent\textit{When useful:} several hypotheses are folded definitions / conjunctions to open at once\par

The hands-free version of mac-h.  Instead of naming each definition and each conjunction by hand -- (mac-h 'IS-METRIC-SPACE A1), split, (mac-h 'is-metric A2), split, ... -- (mac-h*) keeps unfolding any assumption whose head is a defined predicate and splitting any AND assumption, restarting until a full pass changes nothing.  It records as a SINGLE step, so the proof (and its PDF) is far shorter.  It adds no kernel rule: every unfold is the same kernel-checked mac-h, every split the same ai -- just driven to a fixpoint.  Ask (what-now) to see which hypotheses it would touch first.

\par\smallskip{\raggedright\noindent\texttt{(grind)}\quad\tactkind{composite}\index{grind@\texttt{grind} (tactic)}\par}\nobreak
\noindent Saturate the no-choice moves at the focus: decompose the goal connective (di) and break open the hypotheses (mac-h*), repeatedly, until neither fires.  The whole introduce-everything prefix in one step.\par
\noindent\textit{Uses:} \texttt{and-elim} \texttt{and-intro} \texttt{forall-intro} \texttt{forsome-elim} \texttt{iff-elim} \texttt{iff-intro} \texttt{implies-intro} \texttt{macete-hyp} \texttt{not-elim} \texttt{not-intro} \texttt{or-elim}.\par
\noindent\textit{When useful:} right after sp, or any time the focus has connective/definitional structure -- normalize first\par

The deterministic normalizer.  It keeps applying di -- which strips a leading FORALL/IMPLIES/AND, introducing the bound variables and moving antecedents into your hypotheses -- and mac-h*, which unfolds defined predicates and splits conjunctions in the hypotheses, until the goal head is no longer a connective and no hypothesis is foldable.  None of these moves involves CHOOSING a lemma, so there is never anything to reconsider: grind just puts the focus into normal form (e.g. a metric-law goal becomes `(d s)(x,y) = (d s)(y,x)' with the metric's defining properties sitting unfolded in context).  It collapses the di...di mac-h* prefix that bloats proofs into a single step, and adds no kernel rule.

\par\smallskip{\raggedright\noindent\texttt{(scout [depth [branch [nodes]]])}\quad\tactkind{meta}\index{scout@\texttt{scout} (tactic)}\par}\nobreak
\noindent Speculatively search to a bounded depth across INDEPENDENT scratch branches; RETURNS a nested list (number-of-branches (d b) goal best-partials closing-branches) rather than printing.  Never touches the live proof.\par
\noindent\textit{Uses:} every one of the 64 kernel operations.\par
\noindent\textit{When useful:} you are stuck and want the machine to TRY move-sequences (returns data)\par

The copilot's deep lane.  Where (what-now)/(tt) suggest one move, scout actually TRIES sequences: it clones the focus into a fresh deduction graph per branch (so backtracking is free -- a dead branch is just discarded), and does a BEST-FIRST search (frontier ordered by open subgoals left, ties broken deeper-first so it dives toward a closure) whose alphabet is (grind), the closers (ass/rfl/crs/arith), the top rewrite-index (mac) rules, the parameterless backchain (bc*) lemmas, and the INST LANE -- (inst+ <universal hyp> <context-typed term>), instantiating a `forall' hypothesis at a term the context already types and detaching the guard.  It RETURNS the result as data -- (examined-count (depth branch) goal best-partials closing-branches), where best-partials is ((open-goals-left (form ...)) ...) most-reduced first and closing-branches is ((form ...) ...) shortest first -- so you can pick it apart programmatically; use (scout-show ...) for the readable REPL report.  Every move is the real, kernel-checked tactic, so a closing branch is a genuine proof when you adopt it with (scout-run k).  A CITATION GUARD suppresses circular closures: a branch that discharges the goal by citing a library theorem that is the goal -- DIRECTLY (statement alpha-equal to the goal, `P proved by P') or TRANSITIVELY (the cited theorem was itself proven using the goal, i.e. a name alpha-equal to the goal is in its proof-debt closure -- the `compact=>complete by compact=>bongo + bongo=>complete' trap) -- is dropped from closing-branches, the theorem name recorded in *scout-citations*, and scout-show notes `the goal is already theorem X'.  (Loops built purely from independent ASSERTED facts have no live dependency edge to detect; that is warrant/debt hygiene, not scout's to check.)  Bounds: depth (default 6), per-node fan-out branch (default 3), total nodes (default 600).  El cheapo: it does not reason about which move is wise (only the witness terms are guessed, from what the context already types -- a bad guess just dies on the clone); it brute-forces a small tree and you eyeball the survivors.  The inst lane closes the metric laws past grind; goals that need a witness scout can't type (a fresh existential, a constructed term) still surface under best-partials.

\par\smallskip{\raggedright\noindent\texttt{(scout-show [depth [branch [nodes]]])}\quad\tactkind{meta}\index{scout-show@\texttt{scout-show} (tactic)}\par}\nobreak
\noindent Like (scout), but PRINTS the human-readable report (examined count, goal, closing branches or best partials) to the REPL.  Returns the same nested list (scout) does.\par
\noindent\textit{Uses:} every one of the 64 kernel operations.\par
\noindent\textit{When useful:} same as scout but you want the readable report at the REPL\par

The eyeball version of scout: same search, same return value, but it also prints the numbered CLOSING branches (or, when none close, the best partials with how many goals each leaves open).  Use scout-show interactively, (scout) when you want to consume the result as data.

\par\smallskip{\raggedright\noindent\texttt{(scout-run k)}\quad\tactkind{composite}\index{scout-run@\texttt{scout-run} (tactic)}\par}\nobreak
\noindent Adopt closing branch k from the last (scout)/(scout-show) onto the live proof, running its tactics for real (they record and display normally).\par
\noindent\textit{Uses:} no kernel operation: it records no inference.\par
\noindent\textit{When useful:} scout found a closing branch you want to adopt onto the live proof\par

After scout finds CLOSING branches [1], [2], ..., (scout-run k) replays branch k's tactic forms through the real tactics on your live *ps*, from the same focus scout cloned -- so the proof advances and the steps are recorded for the script / PDF exactly as if you had typed them.  Works after either (scout) or (scout-show); both stash the closing branches.

\par\smallskip{\raggedright\noindent\texttt{(prep 'ineq)}\index{prep@\texttt{prep} (tactic)}\par}\nobreak
\noindent Why will a tactic not fire here, and what must be done first?  Runs the tactic's OWN preconditions one at a time, marks each ok / PREP / STOP, names the library lemma that repairs each unmet one, and prints a PLAN it has checked by running it on a scratch clone.  READ-ONLY: the live proof is never touched.  (prep 'ineq) is the only method so far.\par
\noindent\textit{Uses:} no kernel operation: it records no inference.\par

A tactic reports failure as a boolean, so every unmet precondition comes out as the same \#f and the same warning -- (ineq) says `goal not a linear-RR consequence of the named assumptions' whether your goal merely needs a di, or an atom needs a typing fact the library already proves, or the goal is simply false.  prep runs the same predicates separately and tells you WHICH failed.  For ineq the obligations are: the goal must BE an order relation (repair: di, counted by measuring, since one di consumes a typed FORALL and its guard together); every atom must be certified in RR by a LITERAL (IN t RR) scan -- (IN k NN) does not count, the oracle does no subtype reasoning (repair: a coercion lemma, e.g. nn-in-rr); some assumption must be order-shaped, since a fact like not(k=0) is invisible to Farkas (repair: a bridge lemma, e.g. nn-pos-of-nonzero); and the goal must actually follow, which only Fourier-Motzkin decides -- a STOP there means no prepping will help, which is the useful answer.  The repair search is by SHAPE over the whole library and is verified by SPECULATION: a candidate counts only if ineq then fires, not merely because it landed something order-shaped.  Every repair that works is reported, not just the first -- they come out alphabetical, which is no order of merit, and if the goal is itself a library theorem then citing IT is one of the closers.  Worked case: |- forall k in NN. \textasciitilde{}(k=0) => k < 2k.  prep returns (di) (di) (fact 'nn-in-rr 'k) (fact 'nn-pos-of-nonzero 'k) (ineq 4) -- six steps, where the hand-built cut/crs route through k = k+0 < k+k = 2k took thirteen and billed two extra transitivity lemmas.  A tactic joins the table with (prep-method! 'FUBA proc); crs and ass are the obvious next two.

\par\smallskip{\raggedright\noindent\texttt{(subst '(= s t))  |  (subst '(== s t))}\quad\tactkind{rule}\index{subst@\texttt{subst} (tactic)}\par}\nobreak
\noindent Rewrite s -> t throughout the goal, using an equation s = t -- or a quasi-equation s == t -- that is in context, in EITHER orientation (Leibniz substitution).  Reaches every position, operator (head) position included; a binder whose variable the equation mentions is not entered.\par
\noindent\textit{Uses:} \texttt{eq-subst}.\par
\noindent\textit{When useful:} you have an equation s = t (or a quasi-equation s == t) in context and want to rewrite the goal by it -- wherever the term occurs, operator position included\par

Use an equation among your hypotheses to replace s by t everywhere in the goal.  BOTH equalities license it: `=' is VNB's partial equality, and `s = t' already entails s = s and t = t, so t is defined wherever s was and the rewrite never replaces a defined term by an undefined one; `==' is quasi-equality (same definedness, equal where defined), which is a congruence and so substitutes exactly as `=' does -- which you need, since the partial-op recursion and bridge facts are stated with `=='.  The head you pass need not match the head in context, and neither need the orientation: all four of (= s t), (= t s), (== s t), (== t s) are searched for, and the goal is always rewritten s -> t.  LIMIT: the rewrite walk reaches argument positions only, so it is a silent no-op on a term in OPERATOR position -- (subst '(= (VADD md) ...)) will not touch the goal ((VADD md) x y).  Structure accessors are almost always in operator position; use the equation as a macete there (mac on the goal, mac-h on a hypothesis).  (Technically: Leibniz substitution from an in-context equality or quasi-equality.)

\par\smallskip{\raggedright\noindent\texttt{(beta)}\quad\tactkind{rule}\index{beta@\texttt{beta} (tactic)}\par}\nobreak
\noindent Beta-reduce a functoid application in the goal.\par
\noindent\textit{Uses:} \texttt{functoid-beta}.\par
\noindent\textit{When useful:} the goal has a functoid applied to an argument\par

Simplify a function-expression applied to an argument by substituting the argument into its body.  (Technically: beta-reduction of a functoid application in the goal.)

\par\smallskip{\raggedright\noindent\texttt{(nth-r)}\quad\tactkind{rule}\index{nth-r@\texttt{nth-r} (tactic)}\par}\nobreak
\noindent Reduce an NTH applied to a literal LIST in the goal.\par
\noindent\textit{Uses:} \texttt{nth-reduce}.\par
\noindent\textit{When useful:} the goal has NTH of a literal LIST\par

Simplify `the k-th entry of an explicit list [a1, a2, ...]' to that entry.  (Technically: reduce NTH applied to a literal LIST.)

\par\smallskip{\raggedright\noindent\texttt{(len-r)}\quad\tactkind{rule}\index{len-r@\texttt{len-r} (tactic)}\par}\nobreak
\noindent Reduce a LENGTH applied to a literal LIST in the goal.\par
\noindent\textit{Uses:} \texttt{length-reduce}.\par
\noindent\textit{When useful:} the goal has LENGTH of a literal LIST\par

Simplify `the length of an explicit list [a1, ..., an]' to the count n.  Structural (the dual of nth-r): the spine of a list literal is total, so its length is n regardless of whether the entries are defined.

\par\smallskip{\raggedright\noindent\texttt{(rfl)}\quad\tactkind{rule}\index{rfl@\texttt{rfl} (tactic)}\par}\nobreak
\noindent Close a reflexive equality goal (t = t).\par
\noindent\textit{Uses:} \texttt{reflexivity}.\par
\noindent\textit{When useful:} the goal is `t = t' with t manifestly defined\par

Close a goal `t = t' -- a thing equals itself.  (Technically: reflexivity of equality.)

\par\smallskip{\raggedright\noindent\texttt{(qrfl)}\quad\tactkind{rule}\index{qrfl@\texttt{qrfl} (tactic)}\par}\nobreak
\noindent Quasi-reflexivity: close t = t under the partial-equality definedness reading.\par
\noindent\textit{Uses:} \texttt{quasi-reflexivity}.\par
\noindent\textit{When useful:} the goal is `t = t' but t might be UNDEFINED (partial =)\par

Close `t = t' under the partial-equality reading, where asserting t = t also asserts that t is DEFINED.  Use this rather than rfl when t might be undefined.  (Technically: quasi-reflexivity; see the partial-equality convention, where `t = t' is the definedness predicate.)


\subsection*{Propositional logic}

\par\smallskip{\raggedright\noindent\texttt{(prop)}\index{prop@\texttt{prop} (tactic)}\par}\nobreak
\noindent Decide the goal by PROPOSITIONAL logic from the context, and close it if it follows.  Every non-connective formula -- `x in a', an equation, a whole `forall(...)' -- is one opaque atom; AND/OR/NOT/IMPLIES/IFF are read.  On a goal that does NOT follow it prints the countermodel (which atom must be true, which false) and leaves the proof untouched.  Adds no trust: it decides semantically, then discharges through di/ai/oi/ass/use-em/have!/detach!, so the bill is unchanged and the recorded script is the ordinary step-by-step proof.\par
\noindent\textit{Uses:} \texttt{and-elim} \texttt{and-intro} \texttt{arith-forsome} \texttt{arith-ground} \texttt{arith-simplify} \texttt{assumption} \texttt{cartesian-decompose} \texttt{cut} \texttt{detach} \texttt{eq-subst} \texttt{forall-elim} \texttt{forall-intro} \texttt{forsome-elim} \texttt{iff-elim} \texttt{iff-intro} \texttt{implies-intro} \texttt{macete} \texttt{not-elim} \texttt{not-intro} \texttt{or-elim} \texttt{or-intro-left} \texttt{or-intro-right} \texttt{proof-by-contradiction} \texttt{theorem-assumption} \texttt{tuple-equality-decompose}.\par
\noindent\textit{When useful:} the goal follows from the hypotheses by AND/OR/NOT/IMPLIES/IFF alone -- no quantifier or equality reasoning needed\par

Close a goal that follows from the hypotheses by pure propositional reasoning -- the leaves where you can SEE the answer and still have to pick between (oi-l)(ass), (oi-r) plus a conjunction split, and (ai) on a negation.  It treats each atomic statement as a black box, so it will not instantiate a quantifier or reason about equality: if the goal needs an instance, land the instance first (fact / inst+) and run it again.  When it declines it names a falsifying assignment, which usually tells you exactly which hypothesis is missing.  (Technically: three-valued evaluation over the atoms decides the entailment; the proof is then replayed through the kernel rules, case-splitting with excluded middle where a disjunctive goal needs it.)


\subsection*{Arithmetic \& ring oracles}

\par\smallskip{\raggedright\noindent\texttt{(arith)}\quad\tactkind{oracle}\index{arith@\texttt{arith} (tactic)}\par}\nobreak
\noindent Discharge a ground arithmetic goal by evaluation.\par
\noindent\textit{Uses:} \texttt{arith-forsome} \texttt{arith-ground} \texttt{arith-simplify} \texttt{cartesian-decompose} \texttt{eq-subst} \texttt{macete} \texttt{tuple-equality-decompose}.\par
\noindent\textit{When useful:} the goal is ground arithmetic -- no variables\par

Close a goal that is a concrete numerical fact with no variables -- e.g. 2 + 3 = 5, or 7 in NN -- by just computing it.  (Technically: decision by ground arithmetic evaluation.)

\par\smallskip{\raggedright\noindent\texttt{(rs)}\quad\tactkind{oracle}\index{rs@\texttt{rs} (tactic)}\par}\nobreak
\noindent Ring-simplify the goal (normal form over the ambient ring).\par
\noindent\textit{Uses:} \texttt{ring-simplify}.\par
\noindent\textit{When useful:} the goal is a (possibly non-commutative) ring-word identity\par

Simplify the goal to a normal form in the ambient ring (which need not be commutative).  The non-commutative companion of crs.  (Technically: ring-simplify to a canonical word form.)

\par\smallskip{\raggedright\noindent\texttt{(crs)}\quad\tactkind{oracle}\index{crs@\texttt{crs} (tactic)}\par}\nobreak
\noindent Commutative-ring decision procedure: prove a polynomial identity over ZZ[generators].  Expands literal powers, so (x+y)\textasciicircum{}2 = ... closes directly.\par
\noindent\textit{Uses:} \texttt{comm-ring-simplify}.\par
\noindent\textit{When useful:} the goal is a commutative-ring identity, e.g. (x+y)\textasciicircum{}2 = x\textasciicircum{}2+2xy+y\textasciicircum{}2\par

Prove a polynomial identity that holds in EVERY commutative ring -- e.g. (x+y)\textasciicircum{}2 = x\textasciicircum{}2 + 2xy + y\textasciicircum{}2 -- by reducing both sides to a canonical sum-of-monomials form and checking they agree.  A genuine decision procedure: a true commutative-ring identity closes, a non-identity is refused.  Literal powers are expanded for you.  (Technically: normal form over the free commutative ring ZZ[generators].)

\par\smallskip{\raggedright\noindent\texttt{(simp [target])}\quad\tactkind{oracle}\index{simp@\texttt{simp} (tactic)}\par}\nobreak
\noindent Rewrite a commutative-ring SUBTERM of the goal to canonical form, IN PLACE (e.g. (x+y)\textasciicircum{}2 inside a larger goal becomes x\textasciicircum{}2 + 2*x*y + y\textasciicircum{}2).  Works on BOTH surfaces: concrete number domains (+ * - \textasciicircum{} over NN/ZZ/QQ/RR/CC) and a generic ring s ((ADD s)/(MUL s)/(NEG s), carrier (CARR s)).  No arg = outermost ring subterm; ``term'' targets a specific one.  Sound by cut + crs + eq-subst (no new kernel rule); needs the subterm's generators typed in context (true post-di), else refuses and names them.\par
\noindent\textit{Uses:} \texttt{comm-ring-simplify} \texttt{cut} \texttt{eq-subst}.\par
\noindent\textit{When useful:} a ring SUBTERM of a larger goal should be put in normal form in place\par
\par\smallskip{\raggedright\noindent\texttt{(supply)}\index{supply@\texttt{supply} (tactic)}\par}\nobreak
\noindent Close an inequality goal the way a hand proof would: beta-reduce any applied lambda, land the typing certificates and the standard bounds the linear oracle cannot see -- including the TRIANGLE inequality at the summands of an abs of a sum -- then call ineq over every premise.  Rehearses the whole sequence on a throwaway copy and does NOTHING unless it closes; on a miss it prints the sequence it tried.  Adds no trust: it drives lam-b, fact and ineq, each of which records itself.\par
\noindent\textit{Uses:} every one of the 64 kernel operations.\par

The committing form of what-now's SUPPLY THE ORACLE lane.  `ineq' reads abs(...), max(...) and (dist(s))(x,y) as opaque atoms and refuses any premise whose atoms are not certified real, so an inequality that is obviously true can fail for want of a typing nobody mentioned.  This lands them.  The one real idea is subadditivity: where the goal bounds abs(u + v), the useful intermediate term is abs(u) + abs(v), and the decomposition is in the term itself -- nothing is searched for.  Because a lambda reduction cannot be undone and can owe an unprovable leaf, the whole sequence is rehearsed before any of it is run.  (Technically: certificate closure to depth 2 over the forward-citation lane, plus the curated bound table, then Fourier-Motzkin.)

\par\smallskip{\raggedright\noindent\texttt{(ineq i1 i2 ...)}\quad\tactkind{oracle}\index{ineq@\texttt{ineq} (tactic)}\par}\nobreak
\noindent Close a linear-inequality goal over RR (<= < > >= = between RR terms) as a consequence of the named assumptions (1-based indices), via the Fourier-Motzkin/Farkas oracle.  Linearizes over + - * and the binplus/binneg/bintimes aliases; every MAXIMAL non-arithmetic subterm is an atom that must be certified in RR.  (Does NOT see through a generic ring's (ADD s)/(MUL s) -- those become opaque atoms.)\par
\noindent\textit{Uses:} \texttt{ineq}.\par
\noindent\textit{When useful:} a LINEAR <= / < over RR that follows from inequalities you can cite\par

Close a LINEAR inequality over the reals that follows from inequalities you cite (by their hypothesis numbers) -- by chaining them, adding them, and scaling by positive constants.  Anything that is not built from + - and multiplication-by-constants is treated as an opaque quantity, so it handles e.g. the triangle inequality where the distances are unknowns.  For NONLINEAR (polynomial) inequalities use sos instead.  (Technically: Fourier-Motzkin / Farkas over the ordered field RR; every atom must be certified real.)

\par\smallskip{\raggedright\noindent\texttt{(sos ``c1'' ``c2'' ...)}\quad\tactkind{oracle}\index{sos@\texttt{sos} (tactic)}\par}\nobreak
\noindent Sum-of-squares closer for a nonstrict polynomial inequality a <= b over RR (the nonlinear companion of (ineq)).  You supply the terms to be SQUARED; it finds the nonnegative coefficients.  (tactics 'sos) for the worked example.\par
\noindent\textit{Uses:} \texttt{sos}.\par
\noindent\textit{When useful:} a NONSTRICT POLYNOMIAL <= over RR (the nonlinear cousin of ineq)\par

A sum-of-squares closer.  To prove  a <= b  it is enough to exhibit b - a  as a sum of squares (each obviously >= 0), reducing the inequality to an algebraic identity.  You hand sos the terms to be SQUARED -- the c\_i, NOT the squares -- and it solves for nonnegative rationals lambda\_i with

\begin{VNBsnip}
      b - a  =  lambda_1 c_1^2 + ... + lambda_n c_n^2
\end{VNBsnip}

(matched coefficient-by-coefficient in crs's commutative-ring normal form), then PRINTS the lambda\_i it found.  You do not supply them.

Worked example.  Goal  forall([x in rr, y in rr], x*y <= x\textasciicircum{}2 + y\textasciicircum{}2).  Here b - a = x\textasciicircum{}2 - x*y + y\textasciicircum{}2 = 1/2(x-y)\textasciicircum{}2 + 1/2 x\textasciicircum{}2 + 1/2 y\textasciicircum{}2,  so the things squared are  x-y, x, y:

\begin{VNBsnip}
      (sos "x - y" "x" "y")
          ;; closes, printing  1/2(x-y)^2 + 1/2 x^2 + 1/2 y^2
\end{VNBsnip}

2*x*y <= x\textasciicircum{}2 + y\textasciicircum{}2 needs only one square:  (sos ``x - y'').  The three-variable a*b+b*c+c*a <= a\textasciicircum{}2+b\textasciicircum{}2+c\textasciicircum{}2  wants  (sos ``a - b'' ``b - c'' ``c - a'').

Notes. - Write the c\_i with the goal's own variable names.  di preserves them, so sos works before OR after di; skip di and sos peels the typed forall([... in rr]) wrapper itself. - Every generator (variable) must be certified in RR -- a goal binder x in rr  or an  (IN g RR)  assumption. - A wrong or insufficient certificate refuses cleanly: nothing is closed, you simply try other squares. - Strict (<) goals are refused -- a square can be 0, so squares alone never force a strict inequality.


\subsection*{Chained reasoning}

\par\smallskip{\raggedright\noindent\texttt{(calc L0 (rel1 L1 [just1]) (rel2 L2 [just2]) ...)}\quad\tactkind{composite}\index{calc@\texttt{calc} (tactic)}\par}\nobreak
\noindent Ground the focus goal (REL L0 Ln) by a CHAIN of intermediaries L0 rel1 L1 rel2 L2 ... reln Ln: proves each link and composes them into the endpoint.  A link no lane can close is LEFT OPEN as a leaf -- the refinement point, and the PSS candidate.  Relations: = == (folded by transitivity), < <= (folded through the co-*-trans lemmas, with = steps riding along), IFF (chained per direction).  RETURNS the links, the open ones, and their PSS candidates as an alist.\par
\noindent\textit{Uses:} every one of the 64 kernel operations.\par
\noindent\textit{When useful:} the goal is (REL L0 Ln) and you can WRITE the chain of intermediaries that gets there -- or you want the one link that will not close isolated as a PSS candidate\par

Write the argument the way you would on paper -- as a chain through intermediate terms -- and let the machine prove each link:

\begin{VNBsnip}
      Goal  k < 2*k   (with k in NN, ~(k=0))
      (calc 'k '(= (+ k 0)) '(< (+ k k)) '(= (* 2 k)))
          ;; i.e.  k = k+0 < k+k = 2*k
\end{VNBsnip}

Each step names a RELATION and the next LINE; calc cuts the link (rel L(i-1) Li), dispatches a lane at it, and on success composes the whole chain into the goal.  It is pure bookkeeping over the trusted tactics -- no kernel rule, no axiom of its own -- so a calc proof bills exactly the lemmas its lanes and composers cite.

The COMPOSER is forced by the relations you used, and there are three families: `cong' for = and ==, folded by eq-trans; `order' for < and <=, folded through the co-lt-trans / co-le-lt-trans / co-eq-lt-trans ... lemmas (a = step rewrites an endpoint of the running relation and rides along, which is why the chain above needs no separate arithmetic); and `iff', which di's the goal into its two directions and chains each with ai/detach!, because VNB cannot quantify over propositions and so has NO first-order iff-trans lemma to fold with.

The JUSTIFICATION of a link is optional and defaults to 'auto -- try the relation's lanes: crs then arith for = / ==; for < / <= the NN->RR bridge then ineq over the order-shaped premises; grind for IFF.  Otherwise pass 'scout (a small scout search), a macete/theorem NAME (mac it, then close), an explicit tactic form like '(fact 'nn-pos-of-nonzero 'k) eval'd at the link's focus, or 'open / \#f to leave the link open ON PURPOSE.

Sanity first: calc ERRORS before touching the proof if the goal's head is not the composed relation, or its LHS is not L0, or its RHS is not Ln.  A link already in the main context is taken as a GIVEN and skipped -- cutting it would be an alpha self-loop that opens no leaf.

What it is FOR is the open links.  A chain whose links are all closed is a proof; a chain with one open link has isolated your obstacle to a single formula, printed with its free variables and their context typing -- which is the PSS candidate to state as a lemma.  Refine by inserting more intermediaries until each link is discoverable.  (See also (prep 'ineq), which answers the other question: why a lane will not fire.)

\par\smallskip{\raggedright\noindent\texttt{(have! CLAIM [THUNK])}\quad\tactkind{composite}\index{have"!@\texttt{have"!} (tactic)}\par}\nobreak
\noindent Assert CLAIM as an intermediate step and carry on with it as a hypothesis: cut CLAIM, discharge the side goal from context (or by THUNK), and stay on the main branch.\par
\noindent\textit{Uses:} \texttt{and-intro} \texttt{arith-forsome} \texttt{arith-ground} \texttt{arith-simplify} \texttt{assumption} \texttt{cartesian-decompose} \texttt{cut} \texttt{detach} \texttt{eq-subst} \texttt{forall-elim} \texttt{forall-intro} \texttt{iff-intro} \texttt{implies-intro} \texttt{macete} \texttt{not-intro} \texttt{theorem-assumption} \texttt{tuple-equality-decompose}.\par
\noindent\textit{When useful:} you want to state an intermediate fact that follows immediately from context and continue with it -- the `we have X' step\par

The `we have X' step.  You state an intermediate fact CLAIM that follows immediately from what is already known; have! cuts it, proves the resulting side goal automatically, and returns you to the main branch with CLAIM now available as a hypothesis.  The automatic discharge (`from-context!') closes the everyday cases -- a conjunction, splitwise; an (IN (a*b) NN) by nn-mul-closed on the factors; a numeral membership by arith; otherwise the goal is already a context assumption up to alpha.  When the side goal needs more than that, pass a THUNK -- any tactic sequence -- as the second argument to discharge it your way.  It ERRORS, never silently no-ops, if CLAIM is already in context up to alpha: the cut would self-loop, opening one child and no main branch.  So a script reads as the argument does -- `we have q0*q0 = 3*(k*k); we have k =/= 0; ...'.  (Technically: composite -- cut then from-context! or THUNK; adds no kernel rule.)


\subsection*{Backchaining with a theorem}

\par\smallskip{\raggedright\noindent\texttt{(ta 'name)}\quad\tactkind{rule}\index{ta@\texttt{ta} (tactic)}\par}\nobreak
\noindent Theorem-assumption: bring the named installed theorem into context as an assumption.\par
\noindent\textit{Uses:} \texttt{theorem-assumption}.\par
\noindent\textit{When useful:} you want a library theorem available as a hypothesis\par

Bring an already-proved theorem into your current hypotheses so you can use it (instantiate it, detach from it, ...).  (Technically: adds the named installed theorem as an assumption.)

\par\smallskip{\raggedright\noindent\texttt{(wbc ['name])}\quad\tactkind{composite}\index{wbc@\texttt{wbc} (tactic)}\par}\nobreak
\noindent Witness-shape backchain: on an `exists v. ...' goal, cite a PSS lemma that MANUFACTURES a witness of that shape, so you can finish with inst+/grind/ew.\par
\noindent\textit{Uses:} \texttt{and-elim} \texttt{arith-forsome} \texttt{arith-ground} \texttt{arith-simplify} \texttt{cut} \texttt{detach} \texttt{forall-elim} \texttt{forsome-elim} \texttt{iff-elim} \texttt{not-elim} \texttt{or-elim} \texttt{theorem-assumption}.\par
\noindent\textit{When useful:} an EXISTENCE goal whose witness a PSS lemma manufactures (diagonalization / block-family)\par

The discovery move for existence proofs.  Faced with `there exists phi with ...', it looks up the produces-witness index for lemmas whose conclusion builds the RIGHT KIND of object -- the same existential-witness shape -- even when the rest of the statement differs (which is why plain bc* misses them: bc* demands the WHOLE conclusion unify).  E.g. on `exists phi. STRICTLY-MONO-NN(phi) and IS-CAUCHY-SEQ(SUBSEQ f phi)' it reaches for `diagonalization' (which makes a STRICTLY-MONO-NN phi from a nested family); on a nested-family goal it reaches for `block-family'.  With no argument it cites the top retrieved producer; (wbc 'name) cites the one you name.  It brings the lemma in (via ta); you then discharge its hypotheses (inst+), skolemize its existential (grind), and supply its witness to your goal (ew).  Circular `headline' lemmas that already conclude your exact goal are filtered out.  (See (witness-producers goal) for the raw retrieval; what-now lists the (wbc 'name) candidates on any existential goal.)

\par\smallskip{\raggedright\noindent\texttt{(bc impl)}\quad\tactkind{rule}\index{bc@\texttt{bc} (tactic)}\par}\nobreak
\noindent Backchain the goal through an (IMPLIES A B) already in context: if the goal matches B, the new goal is A.\par
\noindent\textit{Uses:} \texttt{backchain}.\par
\noindent\textit{When useful:} you have (P => Q) in context and the goal is Q\par

Work backwards through an implication you already have.  If `P implies Q' is among your hypotheses and your goal is Q, this reduces the goal to proving P.  The everyday `to get Q it's enough to show P'.  (Technically: backchaining the goal through an in-context implication whose conclusion matches.)

\par\smallskip{\raggedright\noindent\texttt{(bc* 'thm [((v val)...)] h1 h2 ...)}\quad\tactkind{composite}\index{bc*@\texttt{bc*} (tactic)}\par}\nobreak
\noindent Matching backchain: unify the theorem's conclusion against the goal, then leave its (instantiated) antecedents as subgoals -- optional handlers hk run on the k-th subgoal.\par
\noindent\textit{Uses:} \texttt{assumption} \texttt{backchain} \texttt{cut} \texttt{forall-elim} \texttt{theorem-assumption}.\par
\noindent\textit{When useful:} a library theorem's CONCLUSION matches your goal\par

Apply a known theorem to your goal.  If you have a theorem `if A and B then C' and your goal is its conclusion C -- matching the theorem's variables to yours, and the names of any dummy/bound variables don't matter (`there exists a net N' applies to a goal `there exists a net F') -- this replaces `prove the goal' with `prove A' and `prove B', the theorem's hypotheses.  The everyday `by Theorem X it suffices to show A and B'.  You may attach a handler to each hypothesis to dispatch it; values the match can't determine you supply as ((v val) ...).  (Technically: peels the theorem's leading universals and implications, matches the conclusion -- alpha-aware on bound variables -- and replays ta/inst/cut/bc automatically.)


\subsection*{Lambda, comprehension \& description}

\par\smallskip{\raggedright\noindent\texttt{(lam-t)}\quad\tactkind{rule}\index{lam-t@\texttt{lam-t} (tactic)}\par}\nobreak
\noindent VNB-LAMBDA typing: reduce (IN (VNB-LAMBDA ...) (FUN A B)) to TWO obligations -- the body, and that the domain is a set.\par
\noindent\textit{Uses:} \texttt{lambda-type}.\par
\noindent\textit{When useful:} the goal types a VNB-LAMBDA into FUN(A,B)\par

Show that a function defined by a formula (`x |-> ...') maps A into B -- reduces to showing that, for an arbitrary input in A, the value lies in B, AND that A is a set.  The second is not a formality: carrying the domain in the term stops one lambda from typing into FUN(A,B) for every A, but says nothing about A being a set, and a lambda over a proper class is not a function.  The lambda's declared domain must also be the FUN's, or the rule does not apply at all.  (Technically: VNB-LAMBDA typing into FUN A B.)

\par\smallskip{\raggedright\noindent\texttt{(lam-b)}\quad\tactkind{rule}\index{lam-b@\texttt{lam-b} (tactic)}\par}\nobreak
\noindent VNB-LAMBDA beta: reduce an applied lambda to its substituted body.\par
\noindent\textit{Uses:} \texttt{lambda-beta}.\par
\noindent\textit{When useful:} the goal has a VNB-LAMBDA applied to an argument\par

Simplify a function `x |-> e(x)' applied to an argument a to e(a) -- the body with the argument substituted in.  A lambda carries its DOMAIN, and it is only defined on that domain, so the reduction is licensed only where a is known to be in it: if the membership is neither in the context nor supplied by an enclosing binder, the step still fires but leaves (IN a A) as an extra subgoal.  Land the typing fact BEFORE the lam-b and no subgoal appears.  (Technically: VNB-LAMBDA beta-reduction.)

\par\smallskip{\raggedright\noindent\texttt{(lam-b-h hyp)}\quad\tactkind{rule}\index{lam-b-h@\texttt{lam-b-h} (tactic)}\par}\nobreak
\noindent VNB-LAMBDA beta in a cited ASSUMPTION -- what mac-h is to mac.\par
\noindent\textit{Uses:} \texttt{lambda-beta-hyp}.\par

Beta-reduce an applied lambda inside a hypothesis, in place.  Needed because a `fact' that instantiates a theorem's function variable at a lambda lands the APPLIED lambda in the CONTEXT, where the goal-side `lam-b' cannot reach it: union-of-opens-open at the identity family g := x |-> x lands is-open(md, big-union(i, fam, (x |-> x)(i))).  Without this the proof must detour through a cut beta-equation and a subst.  Cites nothing, so it adds no debt.

\par\smallskip{\raggedright\noindent\texttt{(sep-set)}\quad\tactkind{rule}\index{sep-set@\texttt{sep-set} (tactic)}\par}\nobreak
\noindent Separation sethood: the separation set \{x in A | p\} is a set.\par
\noindent\textit{Uses:} \texttt{sep-sethood}.\par
\noindent\textit{When useful:} the goal asserts a separation set \{ x in A | p \} is a set\par

Show that a set-builder set \{x in A | p(x)\} is genuinely a set.  (Technically: separation sethood -- a subclass of a set is a set.)

\par\smallskip{\raggedright\noindent\texttt{(sep-mi)}\quad\tactkind{rule}\index{sep-mi@\texttt{sep-mi} (tactic)}\par}\nobreak
\noindent Separation membership intro: prove (IN t \{x in A | p\}).\par
\noindent\textit{Uses:} \texttt{sep-mem-intro}.\par
\noindent\textit{When useful:} the goal is membership in a separation set \{ x in A | p \}\par

Prove that a term t belongs to \{x in A | p(x)\} by showing t lies in A and satisfies the condition p.  (Technically: separation-membership introduction.)

\par\smallskip{\raggedright\noindent\texttt{(sep-me hyp)}\quad\tactkind{rule}\index{sep-me@\texttt{sep-me} (tactic)}\par}\nobreak
\noindent Separation membership elim: split a separation-membership assumption into A-membership and the predicate.\par
\noindent\textit{Uses:} \texttt{sep-mem-elim}.\par
\noindent\textit{When useful:} a hypothesis is separation-set membership -- unpack it\par

From a hypothesis that t belongs to \{x in A | p(x)\}, extract the two facts it packs: t lies in A, and p(t) holds.  (Technically: separation-membership elimination.)

\par\smallskip{\raggedright\noindent\texttt{(comp-mi)}\quad\tactkind{rule}\index{comp-mi@\texttt{comp-mi} (tactic)}\par}\nobreak
\noindent Comprehension membership intro.\par
\noindent\textit{Uses:} \texttt{comp-mem-intro}.\par
\noindent\textit{When useful:} the goal is comprehension-set membership\par

Prove that something belongs to a comprehension set by establishing the set's defining condition for it.  (Technically: comprehension-membership introduction.)

\par\smallskip{\raggedright\noindent\texttt{(comp-me hyp)}\quad\tactkind{rule}\index{comp-me@\texttt{comp-me} (tactic)}\par}\nobreak
\noindent Comprehension membership elim.\par
\noindent\textit{Uses:} \texttt{comp-mem-elim}.\par
\noindent\textit{When useful:} a hypothesis is comprehension-set membership\par

From a comprehension-membership hypothesis, extract its defining condition.  (Technically: comprehension-membership elimination.)

\par\smallskip{\raggedright\noindent\texttt{(iota-d term)}\quad\tactkind{rule}\index{iota-d@\texttt{iota-d} (tactic)}\par}\nobreak
\noindent Definite-description: discharge the IOTA uniqueness obligation for `term'.\par
\noindent\textit{Uses:} \texttt{iota-def}.\par
\noindent\textit{When useful:} the goal involves IOTA and you must discharge its uniqueness obligation\par

Justify `the unique x such that p(x)' (a definite description) by proving that exactly one such x exists.  (Technically: IOTA -- discharge the uniqueness obligation for the described term.)

\par\smallskip{\raggedright\noindent\texttt{(bu-set)}\quad\tactkind{rule}\index{bu-set@\texttt{bu-set} (tactic)}\par}\nobreak
\noindent Big-union sethood.\par
\noindent\textit{Uses:} \texttt{big-union-sethood}.\par
\noindent\textit{When useful:} the goal asserts a BIG-UNION is a set\par

Show that a big union (the union of an indexed family of sets) is itself a set.  (Technically: big-union sethood.)

\par\smallskip{\raggedright\noindent\texttt{(bu-mi w)}\quad\tactkind{rule}\index{bu-mi@\texttt{bu-mi} (tactic)}\par}\nobreak
\noindent Big-union membership intro via the index witness w.\par
\noindent\textit{Uses:} \texttt{big-union-mem-intro}.\par
\noindent\textit{When useful:} the goal is BIG-UNION membership -- you have the index witness\par

Prove that an element lies in a big union by exhibiting one index w whose set already contains it.  (Technically: big-union-membership introduction via the index witness.)

\par\smallskip{\raggedright\noindent\texttt{(bu-me hyp)}\quad\tactkind{rule}\index{bu-me@\texttt{bu-me} (tactic)}\par}\nobreak
\noindent Big-union membership elim.\par
\noindent\textit{Uses:} \texttt{big-union-mem-elim}.\par
\noindent\textit{When useful:} a hypothesis is BIG-UNION membership -- obtain its index\par

From a hypothesis that an element lies in a big union, obtain an index whose set contains it (a fresh name for that index).  (Technically: big-union-membership elimination.)


\subsection*{Conditional terms}

\par\smallskip{\raggedright\noindent\texttt{(if-true t)}\quad\tactkind{rule}\index{if-true@\texttt{if-true} (tactic)}\par}\nobreak
\noindent Reduce an (IF p a b) term on the p branch: spawns p as a subgoal; the continuation gains (= (IF p a b) a).\par
\noindent\textit{Uses:} \texttt{if-true}.\par
\noindent\textit{When useful:} the goal has an (IF p a b) term to evaluate on the p branch\par

Evaluate a conditional term `if p then a else b' on the assumption that p holds: you take on p as a side goal, and may then use that the conditional equals a.  (Technically: if-true reduction, spawning p.)

\par\smallskip{\raggedright\noindent\texttt{(if-false t)}\quad\tactkind{rule}\index{if-false@\texttt{if-false} (tactic)}\par}\nobreak
\noindent Reduce an (IF p a b) term on the not-p branch: spawns (NOT p); the continuation gains (= (IF p a b) b).\par
\noindent\textit{Uses:} \texttt{if-false}.\par
\noindent\textit{When useful:} the goal has an (IF p a b) term to evaluate on the not-p branch\par

Evaluate a conditional term `if p then a else b' on the assumption that p fails: you take on `not p' as a side goal, and may then use that the conditional equals b.  (Technically: if-false reduction, spawning not p.)


\subsection*{Navigation, display \& tacticals}

\par\smallskip{\raggedright\noindent\texttt{(focus n)}\index{focus@\texttt{focus} (tactic)}\par}\nobreak
\noindent Switch the focus to the n-th open goal (1-based).\par
\noindent\textit{Uses:} no kernel operation: it records no inference.\par
\par\smallskip{\raggedright\noindent\texttt{(backup-one)}\quad\tactkind{meta}\index{backup-one@\texttt{backup-one} (tactic)}\par}\nobreak
\noindent Take back the last recorded command: restore the goal, the deduction graph and the proof script to the state before it.  Repeat to walk further back; (sp) clears the stack.  `undo' is an alias.\par
\noindent\textit{Uses:} no kernel operation: it records no inference.\par
\noindent\textit{When useful:} the last command was a mistake -- a greedy `di' that ate the induction, a `cut' you did not mean, a branch you want back\par

The one undo in VNB.  It restores the deduction graph -- not merely the focus -- by rolling back the journalled arrow and grounding writes and DROPPING every sequent and inference node posted since, so no node from the abandoned branch survives to be counted as an open leaf and `qed' cannot be handed a phantom obligation.  *proof-script* and the live trace are rewound with it, so the saved script and the printed proof are the proof you actually kept.  It backs up over the LAST RECORDED command: a command that changed nothing was never recorded and is not a step to back up over.  One thing is deliberately not restored -- *fresh-counter*, the eigenvariable source -- so backing up over an `ai' or `ew' that minted `u\_4' and re-running it mints `u\_5'; the proof is the same, the witness has a different name.  Depth is capped at *vnb-undo-depth* (64).

\par\smallskip{\raggedright\noindent\texttt{(undo)}\quad\tactkind{meta}\index{undo@\texttt{undo} (tactic)}\par}\nobreak
\noindent Alias for (backup-one).\par
\noindent\textit{Uses:} no kernel operation: it records no inference.\par
\noindent\textit{When useful:} the last command was a mistake -- alias for backup-one\par
\par\smallskip{\raggedright\noindent\texttt{(show)}\index{show@\texttt{show} (tactic)}\par}\nobreak
\noindent Redisplay the current proof state.\par
\noindent\textit{Uses:} no kernel operation: it records no inference.\par
\par\smallskip{\raggedright\noindent\texttt{(dial notch)}\index{dial@\texttt{dial} (tactic)}\par}\nobreak
\noindent The PRESENTATION RESOLUTION DIAL: how much of the sequent to show.  r0 kernel S-expression, r1 surface syntax (the default, and the last INVERTIBLE notch -- what is printed re-parses to what was printed), r2 guards and typing hypotheses elided, r3 also structures destructured to their declared slot names and deep subterms elided.  (dial) reports the notch, (dial n) sets it, (dial 'auto) reads the current goal and suggests one, (dial 'why) shows the score behind the suggestion.  For ONE FORMULA rather than the proof state, (dial-wff x [n]) -- x a wff, a ``surface string'' or an S-expression, n defaulting to the dial's current notch: (dial-wff (lookup-theorem 'matmul-assoc) 3).\par
\noindent\textit{Uses:} no kernel operation: it records no inference.\par

It is a FILTER, not a set of modes, and two constraints make it one.  The rungs NEST: turning it up only ever SUBTRACTS, so everything readable at a higher notch is readable below it -- checked in the suite, atom by atom.  And the APERTURE IS ON SCREEN: every notch above r1 names itself on the line and counts what it hid, because a lossy render that does not announce its own lossiness is the defect class of `(f)' printing as `f'.  The auto-notch SUGGESTS a starting notch and then STAYS PUT: it is consulted when you ask, never per sequent, so the lens cannot move while you are reading a proof.

\par\smallskip{\raggedright\noindent\texttt{(dial-wff x notch)}\index{dial-wff@\texttt{dial-wff} (tactic)}\par}\nobreak
\noindent Present ONE FORMULA at a resolution notch, instead of the current proof state: x is a wff, a ``surface string'' or an S-expression, and notch defaults to the dial's current setting.  (dial-wff (lookup-theorem 'matmul-assoc) 3).  Same rungs, same aperture line, no proof required.\par
\noindent\textit{Uses:} no kernel operation: it records no inference.\par

`sequent->string' is the display path for a proof STATE, so the dial reaches a formula only while a proof is holding it.  A theorem looked up, a statement being drafted, a hypothesis pasted from somewhere wants the same dial, and this is it.

\par\smallskip{\raggedright\noindent\texttt{(goal-status)}\index{goal-status@\texttt{goal-status} (tactic)}\par}\nobreak
\noindent One-line summary: done / N open goals.\par
\noindent\textit{Uses:} no kernel operation: it records no inference.\par
\par\smallskip{\raggedright\noindent\texttt{(repeat thunk [cap])}\index{repeat@\texttt{repeat} (tactic)}\par}\nobreak
\noindent Run thunk until it stops changing the proof state (LCF REPEAT).\par
\noindent\textit{Uses:} no kernel operation: it records no inference.\par
\par\smallskip{\raggedright\noindent\texttt{(orelse t1 t2 ...)}\index{orelse@\texttt{orelse} (tactic)}\par}\nobreak
\noindent Run thunks in order, stop at the first that makes progress (LCF ORELSE).\par
\noindent\textit{Uses:} no kernel operation: it records no inference.\par
\par\smallskip{\raggedright\noindent\texttt{(quietly thunk)}\index{quietly@\texttt{quietly} (tactic)}\par}\nobreak
\noindent Run thunk with state-dump output suppressed; returns its value.\par
\noindent\textit{Uses:} no kernel operation: it records no inference.\par

\subsection*{Choosing and naming witnesses}

\par\smallskip{\raggedright\noindent\texttt{(choose ex witness [prover])}\index{choose@\texttt{choose} (tactic)}\par}\nobreak
\noindent Choose eps such that FUBA(eps): cut the existential ex = `forsome([eps], FUBA(eps))', discharge it by exhibiting WITNESS (proving FUBA(WITNESS) by PROVER, or from the context), eliminate it, and RETURN the fresh eigenvariable with FUBA of it landed.\par
\noindent\textit{Uses:} \texttt{and-elim} \texttt{and-intro} \texttt{arith-forsome} \texttt{arith-ground} \texttt{arith-simplify} \texttt{assumption} \texttt{cartesian-decompose} \texttt{cut} \texttt{detach} \texttt{eq-subst} \texttt{forall-elim} \texttt{forall-intro} \texttt{forsome-elim} \texttt{forsome-intro} \texttt{iff-elim} \texttt{iff-intro} \texttt{implies-intro} \texttt{macete} \texttt{not-elim} \texttt{not-intro} \texttt{or-elim} \texttt{theorem-assumption} \texttt{tuple-equality-decompose}.\par

The `we may assume without loss of generality that FUBA(eps)' step: a hypothesis `forall([eps in rr], GUBA(eps))' is about to be used at some eps satisfying FUBA, so choose one first and then instantiate the universal at the name this returns (use-at / obtain-at).  You pay for the existence on the spot -- WITNESS is the value and PROVER (a thunk) shows FUBA holds of it; with no PROVER the side goal is closed from the context.  The rest of the proof then depends on FUBA(eps) alone and never on which witness paid for it, which is what makes it read as `without loss of generality'.  If the existential is already in context the cut is skipped and this is plain elimination; a formula that is not an existential is an error.  choose-pos is this tactic with FUBA frozen at pos-rr and witness 1: `let eps > 0 be given' when nothing gives it.  (Technically: composite -- cut / ew / ai and the AND-splitting of what lands; no kernel rule, no choice principle, and the recorded script is the expansion.)

\par\smallskip{\raggedright\noindent\texttt{(minimize! '(v1 ... vk) GUARD MEASURE)}\quad\tactkind{composite}\index{minimize"!@\texttt{minimize"!} (tactic)}\par}\nobreak
\noindent Choose v1..vk satisfying GUARD with the NN-valued MEASURE as small as possible: lands the witnesses and their minimality, and opens the two obligations any minimisation owes (MEASURE lands in NN; GUARD is satisfiable).\par
\noindent\textit{Uses:} \texttt{and-elim} \texttt{and-intro} \texttt{arith-forsome} \texttt{arith-ground} \texttt{arith-simplify} \texttt{assumption} \texttt{cartesian-decompose} \texttt{cut} \texttt{detach} \texttt{eq-subst} \texttt{forall-elim} \texttt{forall-intro} \texttt{forsome-elim} \texttt{forsome-intro} \texttt{iff-elim} \texttt{iff-intro} \texttt{implies-intro} \texttt{macete} \texttt{not-elim} \texttt{not-intro} \texttt{or-elim} \texttt{reflexivity} \texttt{sep-mem-elim} \texttt{sep-mem-intro} \texttt{theorem-assumption} \texttt{tuple-equality-decompose}.\par
\noindent\textit{When useful:} the goal falls to a `least such' / minimal-counterexample argument -- descent proofs (sqrt 2, sqrt 3 irrational), least-degree or least-pivot witnesses\par

The `least such' / minimal-counterexample step, mechanised.  You give three things: the variables to choose, a GUARD formula they must satisfy, and a natural-number-valued MEASURE term to make small.  On the main branch it hands you fresh eigenconstants w1..wk with two hypotheses -- GUARD holds of them, and nothing satisfying GUARD has a strictly smaller MEASURE -- and it opens exactly the two side goals a minimisation genuinely owes: that MEASURE lands in NN wherever GUARD holds (the TYPE goal), and that GUARD holds somewhere (the NONEMPTY goal).  Everything in between -- forming the value set, well-ordering it, unpacking the least witness, and restating minimality in terms of your variables rather than the set -- is done for you.  Every vi must occur in MEASURE (a variable the measure ignores is not one you are minimising over).  Because it quantifies over the vk at the META level, no lambda and no FUN(U,NN) typing obligation ever enters the logic -- it takes a formula and a term, which is what a driver has in hand, and never forms the set U at all.  Returns (list (w1 ... wk) TYPE-node NONEMPTY-node); either node is \#f when that obligation was already in context up to alpha, so test before refocusing.  (Technically: composite -- it drives cut / di / ai / ew / sep-mi / sep-me / fact / inst / detach! / subst / ass / rfl and adds no kernel rule; its one mathematical appeal is `nn-least-element', the well-ordering of NN, proven from ord-well-ordered.  NO choice principle is used: well-ordering returns a MEMBER of a separation set, and sep-me recovers the witness.)

\par\smallskip{\raggedright\noindent\texttt{(obtain LANE)}\quad\tactkind{composite}\index{obtain@\texttt{obtain} (tactic)}\par}\nobreak
\noindent Run LANE -- a thunk whose forward step lands a `there exists' into context -- then eliminate that existential and RETURN the fresh witness's name, read off the proof state rather than guessed.\par
\noindent\textit{Uses:} every one of the 64 kernel operations.\par
\noindent\textit{When useful:} a forward step yields `there exists ...' and you want to name the witness for later use, without guessing the engine's eigenvariable\par

The `let w be such an x' step, with the eigenvariable named for you.  LANE is a zero-argument procedure whose effect is to land some `there exists v. P(v)' among your hypotheses -- e.g. (lambda () (fact 'nn-3-div-square w)).  obtain runs it, spots the existential that newly appeared, eliminates it (introducing a fresh eigenvariable and landing P at that witness), unpacks any conjunction in the body, and hands back the eigenvariable's NAME -- identified by free-variable set-difference (the symbol now in context that was not there before), so the driver never guesses what the engine called it.  This is the robust way to name a descended witness: (define k (obtain (lambda () ...))), then use k.  Returns \#f (with a note) if the lane landed no existential.  (Technically: composite -- a guarded forward discharge followed by forsome-elim (ai) and AND-splitting.  Plain existential elimination, so it owes nothing and uses no choice.)

\par\smallskip{\raggedright\noindent\texttt{(vlet (n1 ...) FORMER)}\quad\tactkind{composite}\index{vlet@\texttt{vlet} (tactic)}\par}\nobreak
\noindent Bind proof-object names from the current state.  FORMER is (match PATTERN) -- bind each name to its slot in a context formula matching PATTERN -- or (choice [v body]) -- eliminate an in-context existential and bind its witness.\par
\noindent\textit{Uses:} \texttt{and-elim} \texttt{cut} \texttt{forsome-elim} \texttt{iff-elim} \texttt{not-elim} \texttt{or-elim}.\par
\noindent\textit{When useful:} you need to NAME a witness or a matched subterm from the proof state, rather than navigate to it by shape\par

A binding form for the pieces of a proof, so a driver NAMES what it needs instead of navigating to it by shape.  Two FORMERs.  (match PATTERN): the listed names ARE the holes (wildcards) in PATTERN; everything else is literal.  vlet searches the context for a formula or subterm matching PATTERN and binds each name to what filled its slot -- pure selection, no proof step, no obligation; a hard error if nothing matches (you named a piece of something absent).  (choice): eliminate the sole existential in context and bind its witness.  (choice v body): present-else-debt -- if `there exists v. body' is already in context up to alpha, eliminate it; otherwise cut it, LEAVE the existence side-goal open as a debt leaf, and eliminate on the main branch -- either way the witness is bound and body[witness] is landed.  vlet's (choice) and `obtain' are the same underlying mechanism -- eliminate an existential, name the witness by free-variable difference -- differing only in what they take: obtain runs a lane and names what it just produced, vlet names an existential already (or, with a body, about to be) in context.  No choice AXIOM is used; a single witness is plain existential-elimination.  (Technically: a define-syntax expanding to (define n ...) over vlet--match / vlet--choice!; composite, no kernel rule beyond forsome-elim and, for present-else-debt, cut.)


